\documentclass[12pt,a4paper]{article}

\usepackage{geometry}
 \geometry{margin=1in}
\usepackage{amsmath}
\usepackage{physics}
 
 \title{Overview and Implementation of the BB84 Distribution Protocol}
 \author{Ashley Kuta, Uyen Nguyen, Shawn McLain}
 
 \begin{document}
 	\maketitle
 	
 	\paragraph{}
 	
 	The development of quantum key distribution (QKD) protocols spurs from a need to identify and resolve weaknesses in public key cryptosystems such as RSA. Modern cryptosystems rely on the fact that certain mathematical problems such as factoring are computationally complex on classical computers. Using Shor's algorithm, a quantum computer can factor a composite positive integer in $\mathcal{O}\left(\left(\log(N)\right)^3\right)$ operations. In contrast, a classical computer can only implement sub-exponential time complexity algorithms to solve the factoring problem. Therefore, with the advent of quantum computers, existing public key cryptosystems would have glaring weaknesses against attacks. As a result, seeking alternative protocols that adhere to quantum mechanical principles would provide information-theoretic secure key distribution based on fundamental physics as opposed to classical computer performance. % Cite textbook for this paragraph
 	
 	\paragraph{}
 	The foundation of QKD traces back to Stephen Wiesner's (then unpublished) \textit{Conjugate Coding} manuscript where it was proposed that encoding information in incompatible bases can be used for secure communications. He noted that wave guides could be used as a communication channel to send quantum information encoded in the polarization of light. Messages can be constructed using either the vertical-horizontal or circular polarization basis. Consider the usual computational (Z) basis: $\left\{\ket{0}, \ket{1}\right\}$. The two basis states can be represented by photons polarized in the horizontal or vertical direction. These polarizations are orthogonal, so a polarimeter can be oriented such that photons from one of the two states will be detected with certainty. Now, consider the X basis: $\left\{\ket{-}, \ket{+}\right\}$. The two basis states can be represented by photons polarized in the left-handed circular or right-handed circular direction. Again, these two polarizations are orthogonal, so measuring in the X basis means a set of photons polarized in this basis can be effectively filtered. However, if the measurement device is set to measure in the X basis while the photons are prepared in the Z basis, or vice-versa, the probability of the measurement device detecting any of the prepared photons is 50 (e.g. \% $\abs{\bra{0}\ket{+}}^2 = \flatfrac{1}{2}$). Wiesner noted that the two bases are non-commuting observables akin to position and momentum. Only measuring the received information using the correct basis can a deterministic result be discerned. By measuring in the wrong basis, not only is the outcome probabilistic, but uncertainty is introduced because the original state is disturbed by the measurement and no further knowledge about it can be recovered, since identical copies of the unknown, original state can't be created (in accordance with the no-cloning theorem). Additionally, the fact that a state can't be recovered makes the detection of eavesdropping possible by the receiver. Hence, the probabilistic nature of incompatible measurement bases provides a foundation for key distribution protocols because they provide a means of encoding and verifying information. The principles outlined in this manuscript lie at the heart of QKD protocols later developed. % Cite Wiesner's manuscript for this paragraph

	\paragraph{}
	The first QDK protocol was developed in 1984 by Bennett and Brassard. Traditional public key cryptosystems require having a pair of keys, one public key that can be shared over a communication channel, and one private key that must be kept secret for the message to remain secure. The algorithm proposed by Bennett and Brassard creates a shared private key between the sender and receiver, akin to private key cryptosystems. However, unlike private key cryptosystems, the BB84 protocol maintains one of the advantages of public key cryptosystems, which is the distribution of secure random keys over an insecure (quantum) communications channel without the need of a secret key shared a priori, as long as an authenticated classical communication channel is available.

	\paragraph{}
	The protocol starts with the sender (Alice) encoding bits in a quantum state (qubits) and transmitting them over a (necessarily insecure) quantum channel. As described above, a physical realization of this is using vertical-horizontal or circular polarizations as a choice of basis each of which each have two orthogonal states in that basis. The receiver (Bob) must choose which basis to measure the incoming qubits without any knowledge of Alice's choice of basis. Bob has a 50\% chance of picking the correct basis. After every message is sent, Alice and Bob publicly compare their choice of basis chosen for each message, and retain the bits where they both chose the same basis. If the public channel were secure and noiseless, then Alice and Bob would have a shared secret key that can be used for one-time pad encryption.

	\paragraph{}
	For an insecure communication channel, an eavesdropper (Eve) can employ and intercept and resend attack. In this attack, Eve measures Alice's message before it reaches Bob. Similar to Bob, Eve must guess which basis the the message was sent in, with a theoretical success rate of 50\%. Since the quantum state collapses after measurement, Eve must prepare another state that gets forwarded to Bob to avoid detection. Since there is a 50\% chance Eve measured and re-prepared a state in the right basis, the probability that Bob would choose the same basis as Alice and not yield the correct measurement from the resent message is 50\%. Therefore, the theoretical quantum bit error rate (QBER) is $50\% \times 50\% = 25\%$ for this attacking scheme. Such an error rate is above the BB84's maximum tolerance level of 11\%. % https://arxiv.org/abs/quant-ph/0003004

	\paragraph{}
	The QBER can be reduced by classical post-processing step known as information reconciliation, which is an error-correction code performed over the communication channel. It's critical that Bob and Alice perform this correction without leaking too much information to Eve, otherwise the resulting key won't be a secret. There are numerous processes that can be used to find the location of a bit error while minimizing information leakage. One such process is the Cascade protocol. In this protocol, Alice and Bob partition their key strings into blocks of equal size (the block size is inversely proportional to the QBER) and publicly compare the parity of each block. If a block has different parity, then Bob (and Eve) knows that at least one bit flip occurred in his block. Bob and Alice can then perform a binary search on the block by comparing the parity of sub-blocks until a single bit flip is isolated. After an initial passthrough, Alice and Bob permute their key string and repeat the process using a different block size. By doing so, new parity errors in the permuted blocks can be identified, lowering the bit error rate further. Note that this process avoids publicly comparing each bit in the key over the communication channel, otherwise Eve would have knowledge of the entire key. % https://link.springer.com/chapter/10.1007/3-540-48285-7_35

 \end{document}